\subsection{THE EXTENSION THEOREM}

Let \(\mathfrak{g} = \mathfrak{h} + \mathfrak{g}'\) be a Lie algebra which is the direct sum of an ideal \(\mathfrak{g}'\) and a subalgebra \(\mathfrak{h}\) , U the enveloping algebra of \(\mathfrak{g}\) and \(\mathrm{U}' \subset \mathrm{U}\) the enveloping algebra of \(\mathfrak{g}'\) . There exists one and only one \(\mathfrak{g}\) -module structure on \(\mathrm{U}'\) such that: (α) for \(x \in \mathfrak{g}'\) and \(u \in \mathrm{U}'\) , \(x_{\mathrm{U}'}u = -ux\) ; (β) for \(x \in \mathfrak{h}\) and \(u \in \mathrm{U}'\) , \(x_{\mathrm{U}'}u = xu - ux\) (the latter element is certainly in \(\mathrm{U}'\) since the inner derivation of U defined by \(x\) leaves \(\mathfrak{g}'\) and hence \(\mathrm{U}'\) stable). For conditions (α) and (β) define uniquely a linear mapping \(x \mapsto x_{\mathrm{U}'}\) of \(\mathfrak{g}\) into \(\mathcal{L}_{\mathrm{K}}(\mathrm{U}')\) . It therefore suffices to verify that \([x, y]_{\mathrm{U}'} = [x_{\mathrm{U}'}, y_{\mathrm{U}'}]\) ; it is only necessary to consider the following cases:

\begin{enumerate}
\def\labelenumi{(\arabic{enumi})}
\tightlist
\item
  \(x \in \mathfrak{g}', y \in \mathfrak{g}'\) : then
\end{enumerate}

\[
[ x, y ] _ {\mathrm{U} ^ {\prime}} u = - u (x y - y x) = \left(x _ {\mathrm{U} ^ {\prime}} y _ {\mathrm{U} ^ {\prime}} - y _ {\mathrm{U} ^ {\prime}} x _ {\mathrm{U} ^ {\prime}}\right) u;
\]

\begin{enumerate}
\def\labelenumi{(\arabic{enumi})}
\setcounter{enumi}{1}
\tightlist
\item
  \(x \in \mathfrak{h}, y \in \mathfrak{g}'\) : then
\end{enumerate}

\[
\begin{array}{r l} {[ x, y ] _ {\mathrm{U}'} u = - u (x y - y x) = x (- u y) - (- u y) x + (x u - u x) y} \\ & {= (x _ {\mathrm{U}'} y _ {\mathrm{U}'} - y _ {\mathrm{U}'} x _ {\mathrm{U}'}) u;} \end{array}
\]

\begin{enumerate}
\def\labelenumi{(\arabic{enumi})}
\setcounter{enumi}{2}
\tightlist
\item
  \(x \in \mathfrak{h}, y \in \mathfrak{h}\) : then \([x, y]_{\mathrm{U}'}\) and \([x_{\mathrm{U}'}, y_{\mathrm{U}'}]\) are two derivations of \(\mathrm{U}'\) whose restrictions to \(\mathfrak{g}'\) coincide with those of \(\operatorname{ad}_{\mathfrak{g}}[x, y]\) and \([\operatorname{ad}_{\mathfrak{g}}x, \operatorname{ad}_{\mathfrak{g}}y]\) ; hence these derivations are equal.
\end{enumerate}

We shall also consider the dual representation \(x \mapsto -^{t}x_{\mathrm{U}'}\) of \(\mathfrak{g}\) on \(\mathrm{U}'^*\) . For \(x \in \mathfrak{g}', -^{t}x_{\mathrm{U}'}\) is the transpose of right multiplication by \(x\) in \(\mathrm{U}'\) ; the corresponding representation of \(\mathrm{U}'\) is therefore the coregular representation of \(\mathrm{U}'\) .

DEFINITION 1. Let g be a Lie algebra, g' a subalgebra of g and \(\rho'\) a representation of g' on V'. A representation \(\rho\) of g on V is called an extension of \(\rho'\) to g if there exists an injective homomorphism of the g'-module V' into the g'-module V. We also say that the g-module V is an extension of the g'-module V'.

If \(\rho'\) is finite-dimensional and \(\mathfrak{g}'\) is a solvable ideal of \(\mathfrak{g}\) , it is necessary for the existence of a finite-dimensional extension that \([\mathfrak{g},\mathfrak{g}']\) be contained in the largest nilpotency ideal of \(\rho'\) ( \(\S 5\) , no. 3, Theorem 1).

THEOREM 1 (Zassenhaus). Let \(\mathfrak{g} = \mathfrak{g}' + \mathfrak{h}\) be a Lie algebra which is the direct sum of an ideal \(\mathfrak{g}'\) and a subalgebra \(\mathfrak{h}\) and \(\rho'\) a finite-dimensional representation of \(\mathfrak{g}'\) whose largest nilpotency ideal contains \([\mathfrak{h},\mathfrak{g}']\) .

\begin{enumerate}
\def\labelenumi{(\alph{enumi})}
\item
  There exists a finite-dimensional extension \(\rho\) of \(\rho'\) to \(\mathfrak{g}\) whose largest nilpotency ideal contains that of \(\rho'\) .
\item
  If for all \(x \in \mathfrak{h}\) the restriction to \(\mathfrak{g}'\) of \(\operatorname{ad}_{\mathfrak{g}}x\) is nilpotent, \(\rho\) can be chosen such that moreover the largest nilpotency ideal of \(\rho\) contains \(\mathfrak{h}\) .
\end{enumerate}

Let \(U'\) be the enveloping algebra of \(g'\) . Suppose that \(U'\) and \(U'^*\) have the \(g\) -module structures defined at the beginning of this no.

Let \(I \subset U^{\prime}\) be the kernel of \(\rho^{\prime}\) (identified with a representation of \(U^{\prime}\) ). It is a two-sided ideal of \(U^{\prime}\) of finite codimension. The subspace \(C(\rho^{\prime})\) of \(U^{\prime*}\) (cf.~no. 1) is orthogonal to I. Let S be the sub-g-module of \(U^{\prime*}\) generated by \(C(\rho^{\prime})\) .

\[
\mathrm{U}^{\prime *} \supset \mathrm{S} \supset \mathrm{C}(\rho')
\]

We now show that S is finite-dimensional over K. Let \(V'\) be the space on which \(\rho'\) operates and \(V' = V_{0}' \supset V_{1}' \supset \cdots \supset V_{d}' = \{0\}\) a Jordan-Hölder series of the \(g'\) -module \(V'\) . Let \(\rho_{i}'\) be the representation of \(g'\) on \(V_{i-1}'/V_{i}'\) derived from \(\rho' (1 \leqslant i \leqslant d)\) . Let \(I' \subset U'\) be the intersection of the kernels of the \(\rho_{i}'\) (identified with representations of \(U'\) ). Then

\[
\mathrm{I} ^ {\prime d} \subset \mathrm{I} \subset \mathrm{I} ^ {\prime}
\]

and \(\mathrm{I}'\cap \mathfrak{g}'\) is the largest nilpotency ideal of \(\rho'\) . By §2, no. 6, Corollary to Proposition 6, \(\mathrm{I}^{\prime d}\) is of finite codimension in \(\mathrm{U}'\) . If \(x\in \mathfrak{h}\) , the derivation \(u\mapsto xu - ux\) of \(\mathrm{U}'\) maps \(\mathfrak{g}'\) into \([\mathfrak{h},\mathfrak{g}']\subset \mathrm{I}'\) , hence \(\mathrm{U}'\) into \(\mathrm{I}'\) and hence \(\mathrm{I}^{\prime d}\) into \(\mathrm{I}^{\prime d}\) . On the other hand, clearly \(\mathrm{I}^{\prime d}\) is a sub- \(\mathfrak{g}'\) -module of \(\mathrm{U}'\) . Hence \(\mathrm{I}^{\prime d}\) is a sub- \(\mathfrak{g}\) -module of \(\mathrm{U}'\) . The orthogonal of \(\mathrm{I}^{\prime d}\) in \(\mathrm{U}^*\) is a finite-dimensional sub- \(\mathfrak{g}\) -module which contains \(\mathrm{C}(\rho')\) and therefore S. This shows that S is finite-dimensional over K. For \(x\in I'\cap \mathfrak{g}', x^d\) is obviously contained in the annihilator of the \(\mathfrak{g}\) -module \(\mathrm{U}' / \mathrm{I}^{\prime d}\) and hence also in the annihilator of the \(\mathfrak{g}\) -module S.

We saw in no. 1 that the \(\mathfrak{g}'\) -module \(\mathbf{V}'\) is isomorphic to a sub- \(\mathfrak{g}'\) -module of a product \((\mathrm{C}(\rho'))^n\) . Hence the \(\mathfrak{g}\) -module \(\mathrm{S}^n\) provides a finite-dimensional extension \(\rho\) of \(\rho'\) to \(\mathfrak{g}\) . Moreover, \(\rho(x)\) is nilpotent for \(x \in \mathrm{I}' \cap \mathfrak{g}'\) ; as \(\mathrm{I}' \cap \mathfrak{g}'\) is an ideal of \(\mathfrak{g}\) (for it contains \([\mathfrak{h},\mathfrak{g}']\) by hypothesis), we see that \(\mathrm{I}'\cap \mathfrak{g}'\) is contained in the largest nilpotency ideal of \(\rho\) . Thus (a) is proved.

Suppose finally that for all \(x \in \mathfrak{h}\) the restriction to \(\mathfrak{g}'\) of \(\mathrm{ad}_{\mathfrak{g}}x\) is nilpotent. As the elements of \(\mathfrak{h}\) operate on \(\mathrm{U}'\) by derivations, there exists for all \(u \in \mathrm{U}'\) and all \(x \in \mathfrak{h}\) an integer \(e\) such that \((x_{\mathfrak{U}'})^e.u = 0\) ; the endomorphisms derived from \(x_{\mathfrak{U}'}\) , on \(\mathrm{U}' / \mathrm{I}^{\prime d}\) and on S (which are finite-dimensional spaces) are therefore nilpotent. Thus \(\rho(x)\) is nilpotent for all \(x \in \mathfrak{h}\) . We have seen earlier that this is also true for \(x \in \mathrm{I}' \cap \mathfrak{g}'\) . As \(\mathrm{I}' \cap \mathfrak{g}'\) is an ideal of \(\mathfrak{g}'\) containing \([\mathfrak{h}, \mathfrak{g}']\) , the sum \(\mathfrak{h} + (\mathrm{I}' \cap \mathfrak{g}')\) is also an ideal of \(\mathfrak{g}\) . Assertion (b) of Theorem 1 then results from the following lemma:

Lemma 1. Let \(\mathfrak{g} = \mathfrak{g}' + \mathfrak{h}\) be a Lie algebra which is the sum of an ideal \(\mathfrak{g}'\) and a subalgebra \(\mathfrak{h}\) . Let \(\sigma\) be a finite-dimensional representation of \(\mathfrak{g}\) . Suppose that \(\sigma(x)\) is nilpotent for all \(x \in \mathfrak{g}'\) and all \(x \in \mathfrak{h}\) . Then \(\sigma(x)\) is nilpotent for all \(x \in \mathfrak{g}\) .

Passing to the quotient by the kernel of \(\sigma\) , \(\sigma\) may be assumed to be faithful. Then \(\mathfrak{g}'\) and \(\mathfrak{h}\) are nilpotent and hence \(\mathfrak{g}\) , which is an extension of a quotient of \(\mathfrak{h}\) by \(\mathfrak{g}'\) , is solvable. Then \(\mathfrak{h}\) and \(\mathfrak{g}'\) are contained in the largest nilpotency ideal of \(\sigma\) (§ 5, no. 3, Corollary 6 to Theorem 1).

For an improvement of Theorem 1, cf.~Exercise 4.

